\subsubsection{1.4.1 本轮修订范围与实现证据}

4.8 版以 2026 年 11 月 11 日可回归的本地系统为核对对象，保留课程要求的六章目录。现行用例以 UC-01～UC-15 为唯一编号基线；旧文档 UC01～UC24 的语义映射详见《旧新用例映射》，不按相同数字误作同一功能。价格、交付、鉴定和人员规则以《系统规约》及 `requirement.txt` 的 BR-01～BR-08 为准；演示政策仍待 BSR 批准。本地启动防护属于 NFR-01/06 的基础设施验收，不新增业务用例。

\BSRKeyPoint{本轮主线}{角色协作与对象授权}{顾客购物车与退出登录已补齐；所有非顾客角色可打卡；经理可查看每日排班和到岗；维修人员有基本档案和完整处理流程；库存、排班和考勤在角色之间自动同步。}

\begin{BSRTable}
\rowcolor{TableHead}\color{white}\bfseries 能力范围 & \color{white}\bfseries 当前状态与后续边界 \\
已实现的本地业务 & 顾客目录、购物车、报价、预约、取货码、逐件交付与归还、异常押金复核、库存调拨与接收、库存流水、维修处理与经理验收、每日排班、上下班打卡、角色权限和审计。数据保存在 SQLite。\\
演示适配器 & 支付和退款只生成演示流水，不发生真实资金划转。定位采用粗略地区边界映射；参考汇率是演示快照，不是实时金融数据。\\
已实现的界面同步 & SSE 推送变更信号，浏览器按当前角色重新查询数据；编辑中的表单暂缓刷新。不同标签页可分别保留角色会话。\\
保留的建设目标 & 真实第三方支付/退款与对账、合格电子签署、可靠通知队列、正式汇率服务、类别树和独立合同明细等仍是目标。租中变更与部分归还已有本地演示路径，不得归为完全未实现。\\
证据等级 & 代码、数据库结构、71 项后端自动测试及隔离数据库浏览器流程检查可证明对应演示路径；013 启动闸门 4 项测试、Windows junction 静态逃逸 GET/HEAD 和本地恢复演练已核验。不能证明真实资金到账、BSR 现状、生产性能、正式认证或云端部署。\\
\end{BSRTable}

第 5 章和第 6 章给出分析类到关系模型的分模块映射，并由《现有目标数据差异》标明 SQLite 现状。概念模型不等于数据库已有同名表；安全检查目前主要保存在工单字段和事件中，实时同步修订号属于内存状态，均不得误建为当前独立业务表。

\subsubsection{1.4.2 版本适用与组织制度边界}

本版供课程评审、开发维护和演示验收使用，采用课程设计需求文档的封面、六章目录和审阅确认形式。学生姓名为万锦融，班级为信管2402。组织背景与现状仍来自 BSR 案例；本轮未开展新的实地访谈，不把界面改进推断成组织现有事实。

所有价格和维修材料费用均按 BSR 挪威市场的 NOK 制定，不能把金额解释为人民币，也不能直接与国内机构价格逐项等同。顾客定位影响参考展示币种，取货门店和有效价目决定正式 NOK 报价。原有实时汇率时效要求作为待接入后的验收目标保留。
